\chapter{Бесконтактные платежи}\label{ch:contactless}
Карта, телефон и~кольцо предъявляются к~терминалу одним жестом
и~используют один радиоинтерфейс, но устройства, ключи и~проверка
держателя у~них разные.

\section{NFC и ISO 14443}\label{sec:contactless-nfc}

Бесконтактная карта общается с~терминалом через NFC (Near Field
Communication). Это радиоинтерфейс ближнего поля на частоте 13,56~МГц,
рассчитанный на короткую дистанцию. Физический и
канальный уровни такого обмена определяет стандарт ISO~14443, который
задаёт типы карт, антиколлизию (процедуру опознания одной карты,
когда в~поле сразу несколько) и транспортный уровень, поверх
которого уже работает платёжное приложение~\cite{iso-14443}.

ISO~14443 определяет два типа инициализации, Type~A и Type~B.
Type~A использует модуляцию ASK~100\% (Amplitude-Shift Keying,
амплитудная манипуляция) и модифицированное кодирование Миллера
при передаче от терминала к~карте, тогда как Type~B использует ASK~10\%
и кодирование NRZ (Non-Return-to-Zero, без возврата к~нулю).
Большинство платёжных карт работают по Type~A,
но терминал для сертификации EMVCo Level~1 (уровень радио- и
низкоуровневого протокола; Level~2~--- ядро, Level~3~--- приложение)
обязан поддерживать оба типа. Поддержка обоих типов~--- одна из первых
точек совместимости, которую проверяют при допуске терминала
к~бесконтактным платежам~\cite{iso-14443,emvco-contactless-book-a}.

У~бесконтактной карты нет собственного источника питания. Поле,
которое создаёт антенна терминала, наводит ток в~антенне карты,
и~этого тока достаточно, чтобы запустить чип и выполнить
криптографические операции. По той же причине бесконтактный
платёж возможен в~кольце, брелоке или стикере, где антенна и~чип
внутри, а терминал сам обеспечит энергию.

Антенна терминала должна создавать поле достаточной силы в~рабочей
зоне до 4~см от поверхности. Металлический корпус,
расположение антенны внутри устройства и толщина защитного стекла
влияют на стабильность связи. Если карта получает недостаточно
энергии, чип может не успеть завершить криптографическую операцию
до выхода из поля, и транзакция сорвётся. Чтобы избежать таких
срывов, ISO~14443 задаёт Frame Waiting Time (FWT)~--- максимальное
время, за которое карта обязана ответить на
команду~\cite{iso-14443,emvco-contactless-book-a}.

\begin{figure}[tbp]
  \centering
  \includefiguresvg[width=\linewidth]{assets/figures/ch08-contactless/contactless-protocol-stack}
  \caption{Стек протоколов бесконтактного платежа и~анатомия
  APDU-команд~\cite{emvco-contactless-book-a,emvco-contactless-book-b,
  emvco-book3,iso-7816-4,iso-7816-5}.}%
  \label{fig:contactless-protocol-stack}
\end{figure}

\section{Бесконтактный EMV: протокол и~ядра}\label{sec:contactless-emv}

Бесконтактная транзакция EMV сжимает контактную схему до короткого
обмена. Спецификации выстраивают последовательность операций так,
чтобы терминал успевал завершить обмен, пока карта ещё находится
в~поле~\cite{emvco-contactless-book-a}.

Поверх транспортного уровня T=CL (contactless transport protocol)
по ISO~14443-4 работает APDU-уровень (Application Protocol Data
Unit) ISO/IEC~7816-4, по которому передаются команды SELECT,
GPO и~другие сообщения платёжного приложения (см.~\autoref{sec:emv-app-selection})~\cite{iso-7816-4}.

Бесконтактный выбор приложения начинается с~команды SELECT к PPSE
\texttt{2PAY.SYS.DDF01}, бесконтактному аналогу PSE
из контактной транзакции. На уровне байтов команда выглядит как
\texttt{00 A4 04 00 0E 32 50 41 59 2E 53 59 53 2E 44 44 46 30 31 00}:
CLA, INS, P1\slash P2, длина и~ASCII-строка имени PPSE.
PPSE возвращает список AID,
терминал выбирает подходящее приложение, затем отправляет GET
PROCESSING OPTIONS (GPO), получает AIP и AFL и считывает только
нужные записи. В~простом сценарии на этом
диалог почти заканчивается, и~карта сразу переходит к~формированию
криптограммы по команде GENERATE~AC, у~которой байт P1 задаёт
тип запрашиваемой криптограммы~--- AAC, TC или ARQC, а~бит
\texttt{0x01} запрашивает подпись CDA~\cite{emvco-contactless-book-a}.

Ниже порога CVM limit (Cardholder Verification Method limit) схема
разрешает операцию без проверки держателя, а~выше запрашивает PIN,
биометрию или иной метод~\cite{visa-tadg}.

\subsection*{Концепция ядра (kernel)}

EMVCo определяет несколько ядер бесконтактного протокола. У~каждой
схемы свой набор правил аутентификации, допустимых сценариев
и~логики офлайн-одобрения, и~EMVCo приводит эти различия к~общему
интерфейсу: терминал говорит с~картой одинаково, а~внутренняя
политика остаётся на~стороне каждой сети~\cite{emvco-contactless-book-a}.

После выбора приложения терминал сопоставляет AID с~нужным ядром
и~передаёт управление соответствующей реализации. Поддержка
новой сети означает новое ядро, новый набор тестов и~новый класс
совместимости. Исторически каждое ядро публиковалось отдельной
Book C-x, но с~2022 года EMVCo выпускает унифицированную
\emph{Contactless Kernel Specification} (Book C-8), которая
сводит общие требования к~ядрам в~один документ и~лицензируется
напрямую у~EMVCo~\cite{emvco-kernel-c8}. Visa, Mastercard и Мир используют разные наборы
бесконтактных правил для терминала, и~детали сертификации разбираются
в~главах о~конкретных системах~\ref{ch:visa},~\ref{ch:mastercard}
и~\ref{ch:mir}.

PPSE возвращает в~FCI список Directory Entry (тег 61); каждая
запись несёт ADF Name (4F), Application Label (50) и~Application
Priority Indicator (87). Combination Selection пересекает AID карты
с~AID-листом терминала и~выбирает запись с~меньшим значением
Priority. На реальной карте Мир--Maestro, которую ВТБ эмитирует
с~2019 года~\cite{mir-vtb-maestro-2019}, терминал в~России
предпочитает приложение Мир с~Priority~01. Само ядро определяется
по~RID~--- первым пяти байтам AID по~ISO/IEC~7816-5~\cite{iso-7816-5};
нумерация ядер закреплена в~EMVCo Book~C, где Kernel~3~--- VCPS Visa,
а~Kernel~2~--- M\slash Chip Mastercard. Сроки действия приложения (теги 5F24\slash 5F25)
проверяются позже, после SELECT~AID и~GPO, на этапе Processing
Restrictions, и~отражаются в~TVR байт~2 (\autoref{tab:tvr-byte2}).

\begin{figure}[tbp]
  \centering
  \includefiguresvg[width=.95\linewidth]{assets/figures/ch08-contactless/kernel-dispatch}
  \caption{Маршрутизация AID к~бесконтактному ядру на примере
  карты Мир--Maestro~\cite{mir-vtb-maestro-2019,iso-7816-5}.}\label{fig:kernel-dispatch}
\end{figure}

\subsection*{qVSDC (Visa)}

qVSDC (quick Visa Smart Debit\slash Credit)~--- бесконтактное ядро Visa.
В~контактном EMV аутентификация и~формирование криптограммы
разделены, а~в~qVSDC они сближены, чтобы карта успевала
ответить в~коротком бесконтактном окне~\cite{visa-core-rules}.

qVSDC~--- бесконтактное ядро карты, тогда как Quick Contactless~---
отдельная конфигурация терминала. Для Quick Contactless Visa прямо требует
режим только онлайн, но для программы Visa Easy Payment Service
(VEPS) публичное руководство допускает как онлайн-, так и~офлайн-
одобрение. Корректнее говорить не \enquote{qVSDC всегда онлайн},
а \enquote{архитектура Visa легко сводится к~быстрой онлайн-проверке,
а конкретный исход зависит от параметров терминала и правил
рынка}~\cite{visa-quick-chip-emv-qvsdc,visa-tadg}.

\subsection*{M/Chip Contactless (Mastercard)}

Бесконтактное ядро Mastercard сочетает динамическую аутентификацию
с~гибким режимом расчёта: ядро отвечает быстро и~при подходящих
параметрах риска одобряет операцию без немедленного обращения
к~хосту~\cite{emvco-contactless-book-c2}.

Офлайн-одобрение работает на~транспорте, парковках и~при малых
покупках, где связь нестабильна или важна минимальная задержка.
Выше пороговой суммы или при иных условиях риска ядро уходит
в~обычную онлайн-проверку.

\subsection*{Мир бесконтактный}

Бесконтактная спецификация Мир реализована в~EMV-апплете MPA
(MIR Payment Application; об~устройстве апплета и~среды
Java~Card\,/\,GlobalPlatform~--- в~разделе~\ref{sec:emv-chip-internals})
и~использует ГОСТ-криптографию
в~аутентификационной цепочке~\cite{nspk-mir-contactless}. Радиоинтерфейс
тот же. Различия лежат на уровне криптографии и~правил сети. Детали MPA,
маршрутизации и~кобейджинга разобраны в~главе~\ref{ch:mir}.

В~таблице~\ref{tab:contactless-kernel-comparison} используются
аббревиатуры из главы~\ref{ch:emv}: fDDA (fast Dynamic Data
Authentication, ускоренная динамическая аутентификация,
\autoref{sec:emv-data-auth}), CVC3 (Card Verification Code~3,
динамический код бесконтактного апплета Mastercard), ARQC
(Authorization Request Cryptogram, криптограмма авторизационного
запроса, \autoref{sec:emv-cryptogram}); ГОСТ~34.10~--- российский
стандарт ЭЦП на~эллиптических кривых; 3DES и~AES~--- симметричные
шифры (см.~главу~\ref{ch:crypto}).

\begin{table}[H]
  \begingroup
  \small
  \setlength{\tabcolsep}{4pt}
  \renewcommand{\arraystretch}{1.2}
  \begin{tabularx}{\textwidth}{
      @{}P{0.20\textwidth}
      Y
      Y
      Y@{}
    }
    \toprule
    \headrow
    \thd{Параметр} &
    \thd{qVSDC (Visa)} &
    \thd{M\slash Chip Contactless (MC)} &
    \thd{Мир бесконтактный} \\
    \midrule
    Аутентификация данных & fDDA & CVC3 или ARQC & ГОСТ 34.10 \\
    Тип криптограммы & ARQC & CVC3 или ARQC & ARQC (ГОСТ) \\
    Офлайн-одобрение & нет (предпочтительно; офлайн
    возможен при fDDA) & да (ограниченно) & да \\
    CVM limit & задаётся схемой по рынку & задаётся схемой по рынку & по правилам НСПК \\
    Криптография & RSA + 3DES (AES) & RSA + 3DES (AES) &
    криптография ГОСТ в~сочетании с~AES \\
    \bottomrule
  \end{tabularx}%
  \endgroup
  \caption{Сравнение основных бесконтактных ядер.}\label{tab:contactless-kernel-comparison}
\end{table}
CVM limit~--- порог, ниже которого бесконтактная транзакция
проходит без запроса PIN или биометрии. Порог задаётся схемой
и~регулятором по~каждому рынку отдельно. Во~Франции и~Германии он
равен 50~EUR, в~Австралии~--- 100~AUD, в~Канаде~--- 250~CAD. В~Великобритании
с~марта 2026 года FCA отменил фиксированный потолок 100~GBP,
и~порог теперь устанавливают сами банки по~собственным
риск-моделям~\cite{visa-core-rules,mc-rules,fca-cvm-limit-2026}. Выше лимита
терминал запрашивает верификацию держателя, а ниже транзакция
проходит в~режиме \enquote{tap and go}. У~продавца с~высоким
средним чеком (ювелирные товары, электроника) большинство бесконтактных
транзакций потребует PIN, тогда как в~общепите и на транспорте большинство
уложится в~лимит.

В~России с~марта 2022 года Visa и Mastercard не проводят внутренние
операции, и~бесконтактные пороги внутрироссийского оборота устанавливаются
ПС~Мир~\cite{nspk-mir-rules}. Числовые значения порогов в~публичную
часть Правил не вынесены и~приводятся в~Стандартах Системы; модель
Offline Операции в~ПС~Мир определена как опциональная~--- Правила
допускают её только для карт с~микропроцессором, для которых эмитент
установил Offline-лимит~\cite{nspk-mir-rules}, поэтому при отсутствии
такого лимита бесконтактные операции по умолчанию уходят на
авторизацию эмитенту в~реальном времени.

\subsection*{Динамические лимиты ридера (DRL)}\label{sec:contactless-drl}

Спецификации Visa и~AmEx позволяют ридеру держать не один набор
бесконтактных лимитов на AID, а~несколько~--- по разным регионам
эмиссии или иным признакам карты. Механизм называется Dynamic
Reader Limit (DRL). У~Visa решающим объектом служит APID
(Application Program Identifier, идентификатор прикладной программы,
тег \texttt{0x9F5A}), который
карта отдаёт в~SELECT response: байты APID кодируют регион
программы и~код страны/валюты, и~ядро терминала по этим
байтам подбирает соответствующий профиль лимитов. У~AmEx похожую
роль играет тег \texttt{0x9F70} с~предзаписанными на~карте DRL Sets.
DRL включают там, где соотношение локальных и~международных
карт сильно неоднородно~--- в~ATM, у~ритейлеров аэропортов,
гостиниц~\cite{visa-tadg}.

\section{Мобильные платежи: SE, HCE и выпуск токена}\label{sec:contactless-mobile}

В~карте и~кольце ключи лежат внутри защищённого чипа и~не~покидают его.
Телефон~--- устройство общего назначения, и~мобильный платёж
выстраивают на~нескольких архитектурах доверия. В~классическом
кошельке используются защищённый элемент (Secure Element, SE)
и~эмуляция карты на стороне хоста (Host Card Emulation, HCE).
Рядом существует другой сценарий, в~котором банковское приложение само
участвует в~ближнем взаимодействии с~терминалом через Bluetooth
Low Energy, BLE~\cite{emvco-mobile,nspk-volna-demo-2025}. У~трёх
сценариев различаются и место хранения ключевого материала, и~способ
получения права действовать от имени карты.

\subsection*{Secure Element (SE)}

Apple Pay использует Secure Element как отдельный защищённый чип
внутри телефона или часов. Он физически изолирован от основного
процессора и операционной системы, хранит ключи и выполняет
криптографию так, чтобы ключевой материал не выходил за пределы
периметра. Скомпрометированная основная ОС не открывает доступ
к SE~\cite{apple-pay-security}.

Архитектурно важнее всего выделенный путь от контроллера NFC
к SE, по которому запрос терминала попадает напрямую в~защищённую среду,
а ответ возвращается тем же маршрутом. Операционная система не видит
ни содержимое ключей, ни сам криптографический обмен.

\subsection*{HCE (Host Card Emulation)}

HCE (Host Card Emulation)~--- программная эмуляция карты, когда
часть платёжной логики вынесена в~операционную систему телефона.
Вместо отдельного чипа платёжная логика опирается на~сервис в~ОС
и~связанный с~ним защищённый раздел; ключи такого кошелька
ограничены по~сроку или числу использований, и~он сильнее зависит
от своевременного обновления через сетевые
сервисы~\cite{android-hce-overview,google-pay-device-tokenization-security}.

SE даёт более жёсткую изоляцию и~меньше зависит от~состояния общей
ОС; HCE платит за~программную гибкость зависимостью от~защищённости
телефона и~непрерывной работы сетевых сервисов.

Две модели сосуществуют по~одной причине~--- контроль над аппаратной
платформой. Apple владеет и~чипом SE,
и контроллером NFC, и правилами доступа к~ним. Исторически
ни один сторонний кошелёк не получал прямого доступа к SE
на iPhone. С~августа 2024 года Apple открыла in-app NFC через
Secure Element для сторонних разработчиков в~семи странах:
США, Великобритания, Канада, Австралия, Новая Зеландия, Япония,
Бразилия. В EU доступ европейским разработчикам кошельков (wallet) открыт
ранее отдельно по Digital Markets Act, через HCE,
без обязательного коммерческого договора с~Apple. Apple Pay получает
сильнейшую аппаратную изоляцию, но за пределами этих юрисдикций
Apple по-прежнему закрывает канал для
остальных~\cite{apple-pay-security,apple-nfc-se-api-2024}. Android
фрагментирован: у~сотен производителей
SE есть только на части устройств; где он есть, доступ контролирует
OEM (Original Equipment Manufacturer, производитель устройства). HCE как программная эмуляция работает на любом устройстве
с~NFC, за счёт чуть более слабой модели доверия, при которой ключи
живут в~ОС вместо отдельного чипа~\cite{android-hce-overview}. Для платёжной системы,
которая не может требовать SE от каждого
производителя, HCE остаётся единственным масштабируемым путём.

\subsection*{Банковское приложение и BLE}

При оплате с iPhone жест ничем не отличается от обычного
бесконтактного. Устройство подносят к~терминалу, и касса ждёт
подтверждения. Внутри при этом может работать другой канал
ближнего обмена, не сводимый к~обычному сценарию NFC.
НСПК описывает технологию \enquote{Волна} как такой канал, который
встраивают в~платёжный сервис или мобильное банковское
приложение через общий SDK, а~терминал должен поддерживать
Bluetooth~\cite{nspk-volna-demo-2025,nspk-volna-sdk-ios}.

SE и~HCE остаются внутри карточной модели бесконтактного обмена,
а~в~BLE-сценарии в~ближнюю фазу операции входит само банковское
приложение и~его связь с~серверной частью банка, и~на~первый план
выходят привязка устройства, подлинность экземпляра приложения,
правила риска и~совместимость с~терминальным парком.
НСПК указывает, что для банков интеграция \enquote{Волны} использует
тот же инструментарий, что и~для NFC~\cite{nspk-volna-integration-2026}.
Доработанные терминалы под \enquote{Волну} принимают одним и~тем же
жестом и~платежи по~СБП, и~операции по~картам
Мир~\cite{nspk-volna-demo-2025}. BLE-сценарий меняет радиоканал
и~состав доверенной среды, но не меняет сам платёжный поток.

\subsection*{Выпуск и~привязка токена: как телефон становится картой}

Если мобильный платёж строится вокруг токена, добавление карты
в~кошелёк запускает конкретную цепочку. Держатель вводит данные
карты в~приложении, кошелёк передаёт их в~токен-сервис карточной
сети (TSP, Token Service Provider; подробнее~\autoref{sec:token-providers}),
который обращается к~эмитенту за подтверждением личности
держателя, а~эмитент при необходимости добавляет собственный шаг
проверки~--- одноразовый код, push-подтверждение или иной механизм.

После одобрения TSP выпускает DPAN (Device PAN), привязанный к
устройству и домену использования, и загружает его вместе
с~ключевым материалом в SE либо в~профиль HCE. С~этого момента
телефон предъявляет себя терминалу как самостоятельный платёжный
инструмент.

HCE-сценарий отличается от SE-сценария тем, что ключевой материал
не хранится в~аппаратном чипе постоянно: вместо мастер-ключа
кошелёк получает ограниченный набор одноразовых ключей (Limited
Use Keys, LUK), которые периодически перевыпускаются токен-сервисом,
а само платёжное приложение исполняется в~обычной среде ОС
с~опорой на программные средства защиты. Шаг ID\&V (Identification
\& Verification, проверка личности держателя при выпуске токена)
может завершиться по-разному: при достаточных данных эмитент выдаёт
\emph{frictionless}-одобрение (без дополнительной проверки),
при недостатке~--- запрашивает дополнительную проверку
(\emph{step-up}~--- эскалация: одноразовый код или
push-подтверждение), а~при подозрениях~--- отказывает в~привязке.

Архитектурно это один и~тот же процесс независимо от бренда
кошелька~\cite{apple-card-provisioning-security,google-pay-device-tokenization-overview}.
Различия между VTS, MDES и TSP НСПК, а~также их операционные
ограничения разобраны в~разделе~\ref{sec:token-providers} и~в~главах
конкретных схем~\ref{ch:visa},~\ref{ch:mastercard} и~\ref{ch:mir}.

\begin{figure}[tbp]
  \centering
  \includefiguresvg[width=\linewidth]{assets/figures/ch08-contactless/contactless-provisioning}
  \caption{Выпуск и~привязка токена при подключении мобильного платежа.}\label{fig:contactless-provisioning}
\end{figure}

\subsection*{CDCVM: биометрия вместо PIN}

При оплате телефоном верификация держателя происходит на самом
устройстве через Face ID, Touch ID, отпечаток пальца или код разблокировки.
Метод называется CDCVM (Consumer Device CVM) и~сообщает эмитенту
только одно~--- локальная проверка на~устройстве прошла
успешно~\cite{emvco-contactless-book-a}. Этого хватает, чтобы
в~типичном мобильном сценарии не отправлять пользователя обратно
к PIN на терминале и учитывать результат при оценке
риска~\cite{visa-tadg,google-pay-device-tokenization-security}.

\subsection*{Что видит эмитент, когда вы платите телефоном}

Авторизационное сообщение мобильного платежа несёт несколько
элементов, которых нет в~обычной бесконтактной оплате картой,
и~эмитент по ним определяет устройство и оценивает
риск~\cite{emvco-payment-tokenisation,visa-tadg}.

Первый признак~--- DE~2, в~котором лежит DPAN, токен вместо настоящего
PAN, выданный TSP при выпуске токена. DPAN выглядит как PAN,
но относится к~отдельному диапазону и представляет конкретное
устройство или домен использования. В DE~2 эмитент видит суррогат,
и~сеть детокенизирует его перед
авторизацией~\cite{emvco-payment-tokenisation}.

Второй признак~--- DE~22, который сообщает, что операция пришла
по бесконтактному каналу. Дополнительные поля EMV несут
криптограмму, счётчики и~признаки, связанные с~устройством
и~проверкой держателя, и~по ним эмитент отличает обычную
бесконтактную карту от мобильного кошелька и~строит антифрод
по наблюдаемым признакам операции~\cite{visa-tadg,mc-tpr}.

Третий признак~--- Token Requestor ID (TRID, идентификатор сервиса,
запросившего выпуск токена; формально разбирается
в~\autoref{sec:token-providers}). Он связывает токен с~кошельком
или сервисом и показывает, через какой канал токен
был выпущен~\cite{emvco-payment-tokenisation}.

На пути к~эмитенту карточная сеть детокенизирует запрос, заменяя DPAN
на FPAN (Funding PAN, исходный номер карты, к~которой привязан
токен), и~вместе с~ним передаются признаки токена. Настоящий PAN
появляется только там, где нужен для авторизационного
решения~\cite{emvco-payment-tokenisation}.

\section{Носимые устройства: кольцо, часы, браслет}\label{sec:contactless-wearables}

В~кольце, часах и~браслете работает тот же базовый принцип, что
и~в~телефоне: антенна NFC, платёжное приложение
и~ядро бесконтактного обмена. По источнику питания носимые
устройства делятся на два класса. Активные устройства (умные часы)
имеют аккумулятор и могут нести собственный защищённый чип.
Пассивные устройства (платёжные кольца и~стикеры) получают питание
только от поля терминала~\cite{emvco-contactless-chip,emvco-contactless-book-a}.

Носимые устройства проходят одобрение уровня~1 по общему
mobile-процессу EMVCo с~дополнительными критериями, учитывающими
форм-фактор: часы или кольцо испытываются по тем же тестам Level~1,
но к ним применяются специфические для wearable требования
к~антенне и~устойчивости связи~\cite{emvco-mobile,emvco-wearable-level-1-approval-process}.
Для пассивных колец это особенно заметно, потому что маленькая антенна и
питание только от поля терминала оставляют меньший запас по
энергии и стабильности обмена.

Умные часы с~собственной ОС получают токен так же, как телефон,
через кошелёк, интеграцию с TSP и~выпуск DPAN. Пассивные кольца и~стикеры обычно
персонализируются при производстве, когда платёжное приложение и ключи
прошиваются в~чип на фабрике, а держатель привязывает кольцо
к~карте через сервис эмитента или через процесс вендора
устройства~\cite{emvco-mobile}.

Ограничения носимых устройств начинаются там, где нужна
верификация держателя. Клавиатуры нет, и~часы обычно полагаются
на контроль ношения на запястье (on-wrist detection): если сняли
и~надели заново, кошелёк требует повторной проверки. Сигнал
полезный, но не эквивалент биометрии, ведь датчик ношения не различает,
кому принадлежит запястье. Пассивные кольца и~стикеры отдельного
метода верификации не дают вовсе, и~их сценарии ограничены суммами
ниже бесконтактного порога.

Без CVM пассивные форм-факторы целиком зависят от порога. Когда
порог повышается, как во многих странах после 2020 года, зона
бесконтактных платежей без проверки расширяется, и~риск при утере
кольца растёт пропорционально. Эмитент компенсирует это собственными
порогами и лимитами по~частоте (velocity) на стороне авторизации.

\begin{figure}[tbp]
  \centering
  \includefiguresvg[width=.8\linewidth]{assets/figures/ch08-contactless/wearable-spectrum}
  \caption{Спектр носимых платёжных устройств: от пассивного кольца до активных часов.}\label{fig:wearable-spectrum}
\end{figure}

\section{Атака ретрансляции и пределы защиты}\label{sec:contactless-relay}

Бесконтактный интерфейс открывает вектор атаки, которого нет
у~контактной карты. Это ретрансляция (relay attack), при которой
злоумышленник ставит одно устройство рядом с~картой жертвы, второе рядом
с~терминалом, и~между ними прокидывает диалог по NFC в~реальном
времени. Терминал видит настоящую карту, карта видит настоящий
терминал, хотя между ними может быть несколько метров и
больше~\cite{francis-practical-relay-attack}.

Криптограмма при этом остаётся корректной, ведь её вычисляет настоящий
чип на основании настоящих данных транзакции, и простая привязка
к~сумме или времени тут бесполезна. Атакующий переносит протокол
на расстояние, не ломая его~\cite{francis-practical-relay-attack}.

Теоретический ответ известен давно: distance bounding, измерение
расстояния по задержке ответа. Идея в~том, чтобы отсекать слишком долгие ответы
и~так выявлять удалённое устройство.
В~массовых терминалах общего назначения distance bounding почти
не~встречается, но Mastercard внедрил собственный вариант~---
Relay-Resistant Protocol (RRP) в~M\slash Chip, с~миллисекундной проверкой
round-trip-time на~этапе обмена данными. Реализация остаётся дорогой,
и~для большинства схем оборудование с~такой точностью не~оправдано
текущим уровнем потерь~\cite{kim-avoine-distance-bounding,mc-paypass-rrp}.

Вместо distance bounding платёжные системы опираются на~многоуровневую
защиту. Первый рубеж~--- бесконтактный порог: ниже CVM limit транзакция
проходит без верификации, но эмитент может дополнительно ограничить
количество и~сумму последовательных операций без PIN. Второй~---
контроль частоты на~стороне эмитента: несколько бесконтактных
операций за короткий интервал или в~нехарактерной для держателя
географии поднимают риск-скоринг. Третий~--- CDCVM в~мобильных
кошельках, где телефон или часы требуют биометрию или код разблокировки
перед каждой операцией, и~ретрансляция одного диалога по~NFC
невозможна без физического участия владельца
устройства~\cite{emvco-contactless-book-a,visa-tadg}.

У~пассивных карт и~колец третьего рубежа нет, и~защита
ложится на~первые два, а~при высоком пороге и~слабых контролях
по~частоте ретрансляция остаётся практичной атакой. Для мобильных
кошельков с~CDCVM она технически возможна, но требует одновременного
доступа к~разблокированному устройству и~ретранслятору~--- сценарий
непрактичный в~большинстве случаев~\cite{francis-practical-relay-attack}.

\begin{figure}[tbp]
  \centering
  \includefiguresvg[width=.75\linewidth]{assets/figures/ch08-contactless/relay-attack}
  \caption{Атака ретрансляции: посредник между картой и терминалом.}\label{fig:relay-attack}
\end{figure}
